%! Introduction to Polynomial Functions — Unit 2, Lesson 2.0 — Student Packet Cover Sheet
\documentclass[10pt]{article}
\usepackage{precalculus-article}
\usepackage{precalculus-boxes}
\usepackage{ltablex}
\keepXColumns

\begin{document}

% Full-bleed plum banner
\begin{tikzpicture}[remember picture, overlay]
  \fill[plum] (current page.north west)
    rectangle ([yshift=-0.9in]current page.north east);
\end{tikzpicture}
\vspace{-0.8in}

\begin{center}
  {\color{white}\textbf{\LARGE Precalculus}}\\[4pt]
  {\color{lilacmid}\large\textbf{Unit 2: Polynomial Functions}}\\[6pt]
  {\color{white}\textbf{\large Lesson 2.0 \quad Introduction to Polynomial Functions}}
\end{center}

\vspace{0.22in}
\namedateperiod
\vspace{0.18in}

\begin{learningtargetbox}
\small
\begin{enumerate}[leftmargin=1.5em, itemsep=5pt]
  \item I can decide whether an expression is a polynomial function, and explain
        why an expression that is not --- a negative exponent, a radical, a
        variable in a denominator or an exponent --- fails the definition.
  \item I can name the degree, leading term, leading coefficient, and constant
        term of a polynomial, and write one in standard form.
  \item I can name the low-degree families --- constant, linear, quadratic, cubic,
        quartic --- and count the turning points on a polynomial graph.
\end{enumerate}
\end{learningtargetbox}

\vspace{0.14in}

\begin{tocbox}
\small
\renewcommand{\arraystretch}{1.5}
\begin{tabularx}{\linewidth}{@{} c l X r @{}}
\rowcolor{lilac}
\textbf{\#} & \textbf{Component} & \textbf{Description} & \textbf{Score} \\
\midrule
1 & Warm-Up        & Spiral review of prerequisite skills            & \blank{1.2cm} \\
2 & Guided Notes   & The polynomial family, its parts, and its shape & \blank{1.2cm} \\
3 & Group Activity & Sorting the polynomial family                   & \blank{1.2cm} \\
4 & Exit Ticket    & Independent check on parts and recognition      & \blank{1.2cm} \\
5 & Homework       & Practice and extension                          & \blank{1.2cm} \\
\midrule
  & \multicolumn{2}{r}{\textbf{Total}} & \blank{1.2cm} \\
\end{tabularx}
\end{tocbox}

\vspace{0.14in}

\begin{remindbox}
\small
This unit is about one family of functions --- the polynomials. Today's job is only
to recognize them and name their parts. Nothing here is hard; it is the vocabulary
every lesson from 2.1 to 2.7 will assume you already own.
\end{remindbox}

\end{document}
